\begin{figure}[h]
	\begin{pcvstack}[boxed, center, space=1em]
		\begin{pchstack}[center, space=1em]
			\procedure[linenumbering]{$\setup(\secparam)$}{
				\pcreturn \egpphgen(\secparam)
			}
			\procedure[linenumbering]{$\kgen(\pkepp)$}{
				\pcparse (\group, \egprime, \ggen, \hgen) = \pkepp \\
				x \sample \Z_{\egprime} \\
				X \leftarrow \ggen^x \\
				pk \leftarrow(\egprime, \ggen, X) \\
				sk \leftarrow(\egprime, \ggen, x) \\
				\pcreturn (sk, pk)
			}
			\procedure[linenumbering]{$\enc(\pkepp, pk, M)$}{
				\pcparse (\group, \egprime, \ggen, \hgen) = \pkepp \\
				y \sample \Z_{\egprime} \\
				Y \leftarrow \ggen^y \\
				T \leftarrow X^y \\
				W \leftarrow T M \\
				\pcreturn (Y, W)
			}
		\end{pchstack}
		\begin{pchstack}[center, space=1em]
			\procedure[linenumbering]{$\dec(\pkepp, sk, ct)$}{
				\pcparse (\group, \egprime, \ggen, \hgen) = \pkepp \\
				(Y, W) \leftarrow ct  \\
				T \leftarrow Y^x \\
				M \leftarrow W T^{-1} \\
				\pcreturn M
			}
			\procedure[linenumbering]{$\kgenlossy(\pkepp)$}{
				\pcparse (\group, \egprime, \ggen, \hgen) = \pkepp \\
				r \sample \bin^{\lambda} \\
				X \leftarrow h^r \\
				pk \leftarrow (\egprime, \ggen, X) \\
				\pcreturn pk
			}
		\end{pchstack}
	\end{pcvstack}
	\caption{A lossy PKE adapted from the ElGamal encryption scheme, $\lelgamal$. }
	\label{fig:folklore_lossyelg}
\end{figure}
